\QHeading{第11题　列竖式计算。}
\par\noindent


\Problem{（1）$2.47+1.85=4.32$}
\par


\MethodTitle{方法一：按题意列竖式}
\Answer{\begin{tikzpicture}[x=1em,y=1em,baseline=(current bounding box.center),draw=answerred,text=answerred]
\node[anchor=base,inner sep=0pt] at (0.000,0) {$2$};
\node[anchor=base,inner sep=0pt] at (0.800,0) {$.$};
\node[anchor=base,inner sep=0pt] at (1.100,0) {$4$};
\node[anchor=base,inner sep=0pt] at (1.900,0) {$7$};
\node[anchor=base,inner sep=0pt] at (0.000,-1.1) {$1$};
\node[anchor=base,inner sep=0pt] at (0.800,-1.1) {$.$};
\node[anchor=base,inner sep=0pt] at (1.100,-1.1) {$8$};
\node[anchor=base,inner sep=0pt] at (1.900,-1.1) {$5$};
\node[anchor=base,inner sep=0pt] at (0.000,-2.3) {$4$};
\node[anchor=base,inner sep=0pt] at (0.800,-2.3) {$.$};
\node[anchor=base,inner sep=0pt] at (1.100,-2.3) {$3$};
\node[anchor=base,inner sep=0pt] at (1.900,-2.3) {$2$};
\node[anchor=base,inner sep=0pt] at (-.7,-1.1) {$+$};
\draw[line width=.5pt] (-1,-1.52)--(2.500,-1.52);
\node[anchor=base,inner sep=0pt,font=\normalsize] at (1.100,.72) {$1$};
\node[anchor=base,inner sep=0pt,font=\normalsize] at (0.000,.72) {$1$};
\end{tikzpicture}}

\MethodBlock{方法二：凑整再减去多加的数}{上方竖式及题目要求的验算仍保留；本法用于核对。把第二个数凑成整数，再减去多加的部分。}{\begin{aligned}[t]&\phantom{={}}2.47+1.85\\
&=2.47+2-0.15\\
&=4.47-0.15\\
&=4.32\end{aligned}}

\MethodBlock{方法三：拆成整数部分和小数部分}{上方竖式及题目要求的验算仍保留；本法用于核对。}{\begin{aligned}[t]&\phantom{={}}2.47+1.85\\
&=2.47+1+0.85\\
&=3.47+0.85\\
&=4.32\end{aligned}}
\Problem{（2）$3.53+2.69=6.22$}

\MethodTitle{方法一：按题意列竖式}
\Answer{\begin{tikzpicture}[x=1em,y=1em,baseline=(current bounding box.center),draw=answerred,text=answerred]
\node[anchor=base,inner sep=0pt] at (0.000,0) {$3$};
\node[anchor=base,inner sep=0pt] at (0.800,0) {$.$};
\node[anchor=base,inner sep=0pt] at (1.100,0) {$5$};
\node[anchor=base,inner sep=0pt] at (1.900,0) {$3$};
\node[anchor=base,inner sep=0pt] at (0.000,-1.1) {$2$};
\node[anchor=base,inner sep=0pt] at (0.800,-1.1) {$.$};
\node[anchor=base,inner sep=0pt] at (1.100,-1.1) {$6$};
\node[anchor=base,inner sep=0pt] at (1.900,-1.1) {$9$};
\node[anchor=base,inner sep=0pt] at (0.000,-2.3) {$6$};
\node[anchor=base,inner sep=0pt] at (0.800,-2.3) {$.$};
\node[anchor=base,inner sep=0pt] at (1.100,-2.3) {$2$};
\node[anchor=base,inner sep=0pt] at (1.900,-2.3) {$2$};
\node[anchor=base,inner sep=0pt] at (-.7,-1.1) {$+$};
\draw[line width=.5pt] (-1,-1.52)--(2.500,-1.52);
\node[anchor=base,inner sep=0pt,font=\normalsize] at (1.100,.72) {$1$};
\node[anchor=base,inner sep=0pt,font=\normalsize] at (0.000,.72) {$1$};
\end{tikzpicture}}

\MethodBlock{方法二：拆成整数部分和小数部分}{上方竖式及题目要求的验算仍保留；本法用于核对。}{\begin{aligned}[t]&\phantom{={}}3.53+2.69\\
&=3.53+2+0.69\\
&=5.53+0.69\\
&=6.22\end{aligned}}

\MethodBlock{方法三：按相同计数单位计算}{上方竖式及题目要求的验算仍保留；本法用于核对。先看成若干个百分之一，按整数加减，最后还原成小数。}{\begin{aligned}[t]&\phantom{={}}3.53+2.69\\
&=(353+269)\div 100\\
&=622\div 100\\
&=6.22\end{aligned}}
\Problem{（3）$8.21-3.57=4.64$}

\MethodTitle{方法一：按题意列竖式}
\Answer{\begin{tikzpicture}[x=1em,y=1em,baseline=(current bounding box.center),draw=answerred,text=answerred]
\node[anchor=base,inner sep=0pt] at (0.000,0) {$8$};
\node[anchor=base,inner sep=0pt] at (0.800,0) {$.$};
\node[anchor=base,inner sep=0pt] at (1.100,0) {$2$};
\node[anchor=base,inner sep=0pt] at (1.900,0) {$1$};
\node[anchor=base,inner sep=0pt] at (0.000,-1.1) {$3$};
\node[anchor=base,inner sep=0pt] at (0.800,-1.1) {$.$};
\node[anchor=base,inner sep=0pt] at (1.100,-1.1) {$5$};
\node[anchor=base,inner sep=0pt] at (1.900,-1.1) {$7$};
\node[anchor=base,inner sep=0pt] at (0.000,-2.3) {$4$};
\node[anchor=base,inner sep=0pt] at (0.800,-2.3) {$.$};
\node[anchor=base,inner sep=0pt] at (1.100,-2.3) {$6$};
\node[anchor=base,inner sep=0pt] at (1.900,-2.3) {$4$};
\node[anchor=base,inner sep=0pt] at (-.7,-1.1) {$-$};
\draw[line width=.5pt] (-1,-1.52)--(2.500,-1.52);
\draw[line width=.35pt] (-0.160,.03)--(0.160,.43);
\node[anchor=base,inner sep=0pt,font=\normalsize] at (0.000,.72) {$7$};
\draw[line width=.35pt] (0.940,.03)--(1.260,.43);
\node[anchor=base,inner sep=0pt,font=\normalsize] at (1.100,.72) {$11$};
\draw[line width=.35pt] (1.740,.03)--(2.060,.43);
\node[anchor=base,inner sep=0pt,font=\normalsize] at (1.900,.72) {$11$};
\end{tikzpicture}}

\MethodBlock{方法二：拆成整数部分和小数部分}{上方竖式及题目要求的验算仍保留；本法用于核对。}{\begin{aligned}[t]&\phantom{={}}8.21-3.57\\
&=8.21-3-0.57\\
&=5.21-0.57\\
&=4.64\end{aligned}}

\MethodBlock{方法三：按相同计数单位计算}{上方竖式及题目要求的验算仍保留；本法用于核对。先看成若干个百分之一，按整数加减，最后还原成小数。}{\begin{aligned}[t]&\phantom{={}}8.21-3.57\\
&=(821-357)\div 100\\
&=464\div 100\\
&=4.64\end{aligned}}
\Problem{（4）$11.56+8.77=20.33$}

\MethodTitle{方法一：按题意列竖式}
\Answer{\begin{tikzpicture}[x=1em,y=1em,baseline=(current bounding box.center),draw=answerred,text=answerred]
\node[anchor=base,inner sep=0pt] at (0.000,0) {$1$};
\node[anchor=base,inner sep=0pt] at (0.800,0) {$1$};
\node[anchor=base,inner sep=0pt] at (1.600,0) {$.$};
\node[anchor=base,inner sep=0pt] at (1.900,0) {$5$};
\node[anchor=base,inner sep=0pt] at (2.700,0) {$6$};
\node[anchor=base,inner sep=0pt] at (0.800,-1.1) {$8$};
\node[anchor=base,inner sep=0pt] at (1.600,-1.1) {$.$};
\node[anchor=base,inner sep=0pt] at (1.900,-1.1) {$7$};
\node[anchor=base,inner sep=0pt] at (2.700,-1.1) {$7$};
\node[anchor=base,inner sep=0pt] at (0.000,-2.3) {$2$};
\node[anchor=base,inner sep=0pt] at (0.800,-2.3) {$0$};
\node[anchor=base,inner sep=0pt] at (1.600,-2.3) {$.$};
\node[anchor=base,inner sep=0pt] at (1.900,-2.3) {$3$};
\node[anchor=base,inner sep=0pt] at (2.700,-2.3) {$3$};
\node[anchor=base,inner sep=0pt] at (-.7,-1.1) {$+$};
\draw[line width=.5pt] (-1,-1.52)--(3.300,-1.52);
\node[anchor=base,inner sep=0pt,font=\normalsize] at (1.900,.72) {$1$};
\node[anchor=base,inner sep=0pt,font=\normalsize] at (0.800,.72) {$1$};
\node[anchor=base,inner sep=0pt,font=\normalsize] at (0.000,.72) {$1$};
\end{tikzpicture}}

\MethodBlock{方法二：凑整再减去多加的数}{上方竖式及题目要求的验算仍保留；本法用于核对。把第二个数凑成整数，再减去多加的部分。}{\begin{aligned}[t]&\phantom{={}}11.56+8.77\\
&=11.56+9-0.23\\
&=20.56-0.23\\
&=20.33\end{aligned}}

\MethodBlock{方法三：拆成整数部分和小数部分}{上方竖式及题目要求的验算仍保留；本法用于核对。}{\begin{aligned}[t]&\phantom{={}}11.56+8.77\\
&=11.56+8+0.77\\
&=19.56+0.77\\
&=20.33\end{aligned}}
\Problem{（5）$0.58+9.64=10.22$}

\MethodTitle{方法一：按题意列竖式}
\Answer{\begin{tikzpicture}[x=1em,y=1em,baseline=(current bounding box.center),draw=answerred,text=answerred]
\node[anchor=base,inner sep=0pt] at (0.800,0) {$0$};
\node[anchor=base,inner sep=0pt] at (1.600,0) {$.$};
\node[anchor=base,inner sep=0pt] at (1.900,0) {$5$};
\node[anchor=base,inner sep=0pt] at (2.700,0) {$8$};
\node[anchor=base,inner sep=0pt] at (0.800,-1.1) {$9$};
\node[anchor=base,inner sep=0pt] at (1.600,-1.1) {$.$};
\node[anchor=base,inner sep=0pt] at (1.900,-1.1) {$6$};
\node[anchor=base,inner sep=0pt] at (2.700,-1.1) {$4$};
\node[anchor=base,inner sep=0pt] at (0.000,-2.3) {$1$};
\node[anchor=base,inner sep=0pt] at (0.800,-2.3) {$0$};
\node[anchor=base,inner sep=0pt] at (1.600,-2.3) {$.$};
\node[anchor=base,inner sep=0pt] at (1.900,-2.3) {$2$};
\node[anchor=base,inner sep=0pt] at (2.700,-2.3) {$2$};
\node[anchor=base,inner sep=0pt] at (-.7,-1.1) {$+$};
\draw[line width=.5pt] (-1,-1.52)--(3.300,-1.52);
\node[anchor=base,inner sep=0pt,font=\normalsize] at (1.900,.72) {$1$};
\node[anchor=base,inner sep=0pt,font=\normalsize] at (0.800,.72) {$1$};
\node[anchor=base,inner sep=0pt,font=\normalsize] at (0.000,.72) {$1$};
\end{tikzpicture}}

\MethodBlock{方法二：拆成整数部分和小数部分}{上方竖式及题目要求的验算仍保留；本法用于核对。}{\begin{aligned}[t]&\phantom{={}}0.58+9.64\\
&=0.58+9+0.64\\
&=9.58+0.64\\
&=10.22\end{aligned}}

\MethodBlock{方法三：按相同计数单位计算}{上方竖式及题目要求的验算仍保留；本法用于核对。先看成若干个百分之一，按整数加减，最后还原成小数。}{\begin{aligned}[t]&\phantom{={}}0.58+9.64\\
&=(58+964)\div 100\\
&=1022\div 100\\
&=10.22\end{aligned}}
\Problem{（6）$10-3.25=6.75$}

\MethodTitle{方法一：按题意列竖式}
\Answer{\begin{tikzpicture}[x=1em,y=1em,baseline=(current bounding box.center),draw=answerred,text=answerred]
\node[anchor=base,inner sep=0pt] at (0.000,0) {$1$};
\node[anchor=base,inner sep=0pt] at (0.800,0) {$0$};
\node[anchor=base,inner sep=0pt] at (1.600,0) {$.$};
\node[anchor=base,inner sep=0pt] at (1.900,0) {$0$};
\node[anchor=base,inner sep=0pt] at (2.700,0) {$0$};
\node[anchor=base,inner sep=0pt] at (0.800,-1.1) {$3$};
\node[anchor=base,inner sep=0pt] at (1.600,-1.1) {$.$};
\node[anchor=base,inner sep=0pt] at (1.900,-1.1) {$2$};
\node[anchor=base,inner sep=0pt] at (2.700,-1.1) {$5$};
\node[anchor=base,inner sep=0pt] at (0.800,-2.3) {$6$};
\node[anchor=base,inner sep=0pt] at (1.600,-2.3) {$.$};
\node[anchor=base,inner sep=0pt] at (1.900,-2.3) {$7$};
\node[anchor=base,inner sep=0pt] at (2.700,-2.3) {$5$};
\node[anchor=base,inner sep=0pt] at (-.7,-1.1) {$-$};
\draw[line width=.5pt] (-1,-1.52)--(3.300,-1.52);
\draw[line width=.35pt] (-0.160,.03)--(0.160,.43);
\node[anchor=base,inner sep=0pt,font=\normalsize] at (0.000,.72) {$0$};
\draw[line width=.35pt] (0.640,.03)--(0.960,.43);
\node[anchor=base,inner sep=0pt,font=\normalsize] at (0.800,.72) {$9$};
\draw[line width=.35pt] (1.740,.03)--(2.060,.43);
\node[anchor=base,inner sep=0pt,font=\normalsize] at (1.900,.72) {$9$};
\draw[line width=.35pt] (2.540,.03)--(2.860,.43);
\node[anchor=base,inner sep=0pt,font=\normalsize] at (2.700,.72) {$10$};
\end{tikzpicture}}

\MethodBlock{方法二：拆成整数部分和小数部分}{上方竖式及题目要求的验算仍保留；本法用于核对。}{\begin{aligned}[t]&\phantom{={}}10-3.25\\
&=10-3-0.25\\
&=7-0.25\\
&=6.75\end{aligned}}

\MethodBlock{方法三：按相同计数单位计算}{上方竖式及题目要求的验算仍保留；本法用于核对。先看成若干个百分之一，按整数加减，最后还原成小数。}{\begin{aligned}[t]&\phantom{={}}10-3.25\\
&=(1000-325)\div 100\\
&=675\div 100\\
&=6.75\end{aligned}}
\par\noindent

\Answer{
\QHeading{计算过程：11. 列竖式计算}

小数点对齐，缺位补零；从右往左逐列计算。下面完整列出进位、退位及各数位的计算过程。

\par\medskip\noindent
\textbf{第1小题：\(2.47+1.85=4.32\)}\par\smallskip
\noindent
\begin{tikzpicture}[x=1em,y=1em,baseline=(current bounding box.center),draw=answerred,text=answerred]
\node[anchor=base,inner sep=0pt] at (0.000,0) {$2$};
\node[anchor=base,inner sep=0pt] at (0.800,0) {$.$};
\node[anchor=base,inner sep=0pt] at (1.100,0) {$4$};
\node[anchor=base,inner sep=0pt] at (1.900,0) {$7$};
\node[anchor=base,inner sep=0pt] at (0.000,-1.1) {$1$};
\node[anchor=base,inner sep=0pt] at (0.800,-1.1) {$.$};
\node[anchor=base,inner sep=0pt] at (1.100,-1.1) {$8$};
\node[anchor=base,inner sep=0pt] at (1.900,-1.1) {$5$};
\node[anchor=base,inner sep=0pt] at (0.000,-2.3) {$4$};
\node[anchor=base,inner sep=0pt] at (0.800,-2.3) {$.$};
\node[anchor=base,inner sep=0pt] at (1.100,-2.3) {$3$};
\node[anchor=base,inner sep=0pt] at (1.900,-2.3) {$2$};
\node[anchor=base,inner sep=0pt] at (-.7,-1.1) {$+$};
\draw[line width=.5pt] (-1,-1.52)--(2.500,-1.52);
\node[anchor=base,inner sep=0pt,font=\normalsize] at (1.100,.72) {$1$};
\node[anchor=base,inner sep=0pt,font=\normalsize] at (0.000,.72) {$1$};
\end{tikzpicture}
\par\smallskip{\normalsize 上方小字表示从后一位进来的数。\par}
\vspace{0pt}
\begin{enumerate}[leftmargin=1.7em,itemsep=5pt,topsep=0pt]
\item 百分位：\(\begin{aligned}[t]&\phantom{={}}7+5\\&=12\end{aligned}\)\quad 本位写2，向前一位进1。
\item 十分位：\(\begin{aligned}[t]&\phantom{={}}4+8+1\\&=12+1\\&=13\end{aligned}\)\quad 本位写3，向前一位进1。
\item 个位：\(\begin{aligned}[t]&\phantom{={}}2+1+1\\&=3+1\\&=4\end{aligned}\)\quad 本位写4。
\end{enumerate}
\par\smallskip
因此，\(\begin{aligned}[t]&\phantom{={}}2.47+1.85\\&=4.32\end{aligned}\)。
\par\medskip

\par\medskip\noindent
\textbf{第2小题：\(3.53+2.69=6.22\)}\par\smallskip
\noindent
\begin{tikzpicture}[x=1em,y=1em,baseline=(current bounding box.center),draw=answerred,text=answerred]
\node[anchor=base,inner sep=0pt] at (0.000,0) {$3$};
\node[anchor=base,inner sep=0pt] at (0.800,0) {$.$};
\node[anchor=base,inner sep=0pt] at (1.100,0) {$5$};
\node[anchor=base,inner sep=0pt] at (1.900,0) {$3$};
\node[anchor=base,inner sep=0pt] at (0.000,-1.1) {$2$};
\node[anchor=base,inner sep=0pt] at (0.800,-1.1) {$.$};
\node[anchor=base,inner sep=0pt] at (1.100,-1.1) {$6$};
\node[anchor=base,inner sep=0pt] at (1.900,-1.1) {$9$};
\node[anchor=base,inner sep=0pt] at (0.000,-2.3) {$6$};
\node[anchor=base,inner sep=0pt] at (0.800,-2.3) {$.$};
\node[anchor=base,inner sep=0pt] at (1.100,-2.3) {$2$};
\node[anchor=base,inner sep=0pt] at (1.900,-2.3) {$2$};
\node[anchor=base,inner sep=0pt] at (-.7,-1.1) {$+$};
\draw[line width=.5pt] (-1,-1.52)--(2.500,-1.52);
\node[anchor=base,inner sep=0pt,font=\normalsize] at (1.100,.72) {$1$};
\node[anchor=base,inner sep=0pt,font=\normalsize] at (0.000,.72) {$1$};
\end{tikzpicture}
\par\smallskip{\normalsize 上方小字表示从后一位进来的数。\par}
\vspace{0pt}
\begin{enumerate}[leftmargin=1.7em,itemsep=5pt,topsep=0pt]
\item 百分位：\(\begin{aligned}[t]&\phantom{={}}3+9\\&=12\end{aligned}\)\quad 本位写2，向前一位进1。
\item 十分位：\(\begin{aligned}[t]&\phantom{={}}5+6+1\\&=11+1\\&=12\end{aligned}\)\quad 本位写2，向前一位进1。
\item 个位：\(\begin{aligned}[t]&\phantom{={}}3+2+1\\&=5+1\\&=6\end{aligned}\)\quad 本位写6。
\end{enumerate}
\par\smallskip
因此，\(\begin{aligned}[t]&\phantom{={}}3.53+2.69\\&=6.22\end{aligned}\)。
\par\medskip

\par\medskip\noindent
\textbf{第3小题：\(8.21-3.57=4.64\)}\par\smallskip
\noindent
\begin{tikzpicture}[x=1em,y=1em,baseline=(current bounding box.center),draw=answerred,text=answerred]
\node[anchor=base,inner sep=0pt] at (0.000,0) {$8$};
\node[anchor=base,inner sep=0pt] at (0.800,0) {$.$};
\node[anchor=base,inner sep=0pt] at (1.100,0) {$2$};
\node[anchor=base,inner sep=0pt] at (1.900,0) {$1$};
\node[anchor=base,inner sep=0pt] at (0.000,-1.1) {$3$};
\node[anchor=base,inner sep=0pt] at (0.800,-1.1) {$.$};
\node[anchor=base,inner sep=0pt] at (1.100,-1.1) {$5$};
\node[anchor=base,inner sep=0pt] at (1.900,-1.1) {$7$};
\node[anchor=base,inner sep=0pt] at (0.000,-2.3) {$4$};
\node[anchor=base,inner sep=0pt] at (0.800,-2.3) {$.$};
\node[anchor=base,inner sep=0pt] at (1.100,-2.3) {$6$};
\node[anchor=base,inner sep=0pt] at (1.900,-2.3) {$4$};
\node[anchor=base,inner sep=0pt] at (-.7,-1.1) {$-$};
\draw[line width=.5pt] (-1,-1.52)--(2.500,-1.52);
\draw[line width=.35pt] (-0.160,.03)--(0.160,.43);
\node[anchor=base,inner sep=0pt,font=\normalsize] at (0.000,.72) {$7$};
\draw[line width=.35pt] (0.940,.03)--(1.260,.43);
\node[anchor=base,inner sep=0pt,font=\normalsize] at (1.100,.72) {$11$};
\draw[line width=.35pt] (1.740,.03)--(2.060,.43);
\node[anchor=base,inner sep=0pt,font=\normalsize] at (1.900,.72) {$11$};
\end{tikzpicture}
\par\smallskip{\normalsize 划线数字是退位前的数，上方小字是退位后参与计算的数。\par}
\vspace{0pt}
\begin{enumerate}[leftmargin=1.7em,itemsep=5pt,topsep=0pt]
\item 十分位退位：\(\begin{aligned}[t]&\phantom{={}}2-1\\&=1\end{aligned}\)\quad 借出1给后一位。
\item 百分位接收退位：\(\begin{aligned}[t]&\phantom{={}}1+10\\&=11\end{aligned}\)\quad 前一位的1变成本位的10。
\item 百分位：\(\begin{aligned}[t]&\phantom{={}}11-7\\&=4\end{aligned}\)\quad 本位写4。
\item 个位退位：\(\begin{aligned}[t]&\phantom{={}}8-1\\&=7\end{aligned}\)\quad 借出1给后一位。
\item 十分位接收退位：\(\begin{aligned}[t]&\phantom{={}}1+10\\&=11\end{aligned}\)\quad 前一位的1变成本位的10。
\item 十分位：\(\begin{aligned}[t]&\phantom{={}}11-5\\&=6\end{aligned}\)\quad 本位写6。
\item 个位：\(\begin{aligned}[t]&\phantom{={}}7-3\\&=4\end{aligned}\)\quad 本位写4。
\end{enumerate}
\par\smallskip
因此，\(\begin{aligned}[t]&\phantom{={}}8.21-3.57\\&=4.64\end{aligned}\)。
\par\medskip

\par\medskip\noindent
\textbf{第4小题：\(11.56+8.77=20.33\)}\par\smallskip
\noindent
\begin{tikzpicture}[x=1em,y=1em,baseline=(current bounding box.center),draw=answerred,text=answerred]
\node[anchor=base,inner sep=0pt] at (0.000,0) {$1$};
\node[anchor=base,inner sep=0pt] at (0.800,0) {$1$};
\node[anchor=base,inner sep=0pt] at (1.600,0) {$.$};
\node[anchor=base,inner sep=0pt] at (1.900,0) {$5$};
\node[anchor=base,inner sep=0pt] at (2.700,0) {$6$};
\node[anchor=base,inner sep=0pt] at (0.800,-1.1) {$8$};
\node[anchor=base,inner sep=0pt] at (1.600,-1.1) {$.$};
\node[anchor=base,inner sep=0pt] at (1.900,-1.1) {$7$};
\node[anchor=base,inner sep=0pt] at (2.700,-1.1) {$7$};
\node[anchor=base,inner sep=0pt] at (0.000,-2.3) {$2$};
\node[anchor=base,inner sep=0pt] at (0.800,-2.3) {$0$};
\node[anchor=base,inner sep=0pt] at (1.600,-2.3) {$.$};
\node[anchor=base,inner sep=0pt] at (1.900,-2.3) {$3$};
\node[anchor=base,inner sep=0pt] at (2.700,-2.3) {$3$};
\node[anchor=base,inner sep=0pt] at (-.7,-1.1) {$+$};
\draw[line width=.5pt] (-1,-1.52)--(3.300,-1.52);
\node[anchor=base,inner sep=0pt,font=\normalsize] at (1.900,.72) {$1$};
\node[anchor=base,inner sep=0pt,font=\normalsize] at (0.800,.72) {$1$};
\node[anchor=base,inner sep=0pt,font=\normalsize] at (0.000,.72) {$1$};
\end{tikzpicture}
\par\smallskip{\normalsize 上方小字表示从后一位进来的数。\par}
\vspace{0pt}
\begin{enumerate}[leftmargin=1.7em,itemsep=5pt,topsep=0pt]
\item 百分位：\(\begin{aligned}[t]&\phantom{={}}6+7\\&=13\end{aligned}\)\quad 本位写3，向前一位进1。
\item 十分位：\(\begin{aligned}[t]&\phantom{={}}5+7+1\\&=12+1\\&=13\end{aligned}\)\quad 本位写3，向前一位进1。
\item 个位：\(\begin{aligned}[t]&\phantom{={}}1+8+1\\&=9+1\\&=10\end{aligned}\)\quad 本位写0，向前一位进1。
\item 十位：\(\begin{aligned}[t]&\phantom{={}}1+0+1\\&=1+1\\&=2\end{aligned}\)\quad 本位写2。
\end{enumerate}
\par\smallskip
因此，\(\begin{aligned}[t]&\phantom{={}}11.56+8.77\\&=20.33\end{aligned}\)。
\par\medskip

\par\medskip\noindent
\textbf{第5小题：\(0.58+9.64=10.22\)}\par\smallskip
\noindent
\begin{tikzpicture}[x=1em,y=1em,baseline=(current bounding box.center),draw=answerred,text=answerred]
\node[anchor=base,inner sep=0pt] at (0.800,0) {$0$};
\node[anchor=base,inner sep=0pt] at (1.600,0) {$.$};
\node[anchor=base,inner sep=0pt] at (1.900,0) {$5$};
\node[anchor=base,inner sep=0pt] at (2.700,0) {$8$};
\node[anchor=base,inner sep=0pt] at (0.800,-1.1) {$9$};
\node[anchor=base,inner sep=0pt] at (1.600,-1.1) {$.$};
\node[anchor=base,inner sep=0pt] at (1.900,-1.1) {$6$};
\node[anchor=base,inner sep=0pt] at (2.700,-1.1) {$4$};
\node[anchor=base,inner sep=0pt] at (0.000,-2.3) {$1$};
\node[anchor=base,inner sep=0pt] at (0.800,-2.3) {$0$};
\node[anchor=base,inner sep=0pt] at (1.600,-2.3) {$.$};
\node[anchor=base,inner sep=0pt] at (1.900,-2.3) {$2$};
\node[anchor=base,inner sep=0pt] at (2.700,-2.3) {$2$};
\node[anchor=base,inner sep=0pt] at (-.7,-1.1) {$+$};
\draw[line width=.5pt] (-1,-1.52)--(3.300,-1.52);
\node[anchor=base,inner sep=0pt,font=\normalsize] at (1.900,.72) {$1$};
\node[anchor=base,inner sep=0pt,font=\normalsize] at (0.800,.72) {$1$};
\node[anchor=base,inner sep=0pt,font=\normalsize] at (0.000,.72) {$1$};
\end{tikzpicture}
\par\smallskip{\normalsize 上方小字表示从后一位进来的数。\par}
\vspace{0pt}
\begin{enumerate}[leftmargin=1.7em,itemsep=5pt,topsep=0pt]
\item 百分位：\(\begin{aligned}[t]&\phantom{={}}8+4\\&=12\end{aligned}\)\quad 本位写2，向前一位进1。
\item 十分位：\(\begin{aligned}[t]&\phantom{={}}5+6+1\\&=11+1\\&=12\end{aligned}\)\quad 本位写2，向前一位进1。
\item 个位：\(\begin{aligned}[t]&\phantom{={}}0+9+1\\&=9+1\\&=10\end{aligned}\)\quad 本位写0，向前一位进1。
\end{enumerate}
最高位再写进上来的1。\par
\par\smallskip
因此，\(\begin{aligned}[t]&\phantom{={}}0.58+9.64\\&=10.22\end{aligned}\)。
\par\medskip

\par\medskip\noindent
\textbf{第6小题：\(10-3.25=6.75\)}\par\smallskip
\noindent
\begin{tikzpicture}[x=1em,y=1em,baseline=(current bounding box.center),draw=answerred,text=answerred]
\node[anchor=base,inner sep=0pt] at (0.000,0) {$1$};
\node[anchor=base,inner sep=0pt] at (0.800,0) {$0$};
\node[anchor=base,inner sep=0pt] at (1.600,0) {$.$};
\node[anchor=base,inner sep=0pt] at (1.900,0) {$0$};
\node[anchor=base,inner sep=0pt] at (2.700,0) {$0$};
\node[anchor=base,inner sep=0pt] at (0.800,-1.1) {$3$};
\node[anchor=base,inner sep=0pt] at (1.600,-1.1) {$.$};
\node[anchor=base,inner sep=0pt] at (1.900,-1.1) {$2$};
\node[anchor=base,inner sep=0pt] at (2.700,-1.1) {$5$};
\node[anchor=base,inner sep=0pt] at (0.800,-2.3) {$6$};
\node[anchor=base,inner sep=0pt] at (1.600,-2.3) {$.$};
\node[anchor=base,inner sep=0pt] at (1.900,-2.3) {$7$};
\node[anchor=base,inner sep=0pt] at (2.700,-2.3) {$5$};
\node[anchor=base,inner sep=0pt] at (-.7,-1.1) {$-$};
\draw[line width=.5pt] (-1,-1.52)--(3.300,-1.52);
\draw[line width=.35pt] (-0.160,.03)--(0.160,.43);
\node[anchor=base,inner sep=0pt,font=\normalsize] at (0.000,.72) {$0$};
\draw[line width=.35pt] (0.640,.03)--(0.960,.43);
\node[anchor=base,inner sep=0pt,font=\normalsize] at (0.800,.72) {$9$};
\draw[line width=.35pt] (1.740,.03)--(2.060,.43);
\node[anchor=base,inner sep=0pt,font=\normalsize] at (1.900,.72) {$9$};
\draw[line width=.35pt] (2.540,.03)--(2.860,.43);
\node[anchor=base,inner sep=0pt,font=\normalsize] at (2.700,.72) {$10$};
\end{tikzpicture}
\par\smallskip{\normalsize 划线数字是退位前的数，上方小字是退位后参与计算的数。\par}
\vspace{0pt}
\begin{enumerate}[leftmargin=1.7em,itemsep=5pt,topsep=0pt]
\item 十位退位：\(\begin{aligned}[t]&\phantom{={}}1-1\\&=0\end{aligned}\)\quad 借出1给后一位。
\item 个位传递退位：\(\begin{aligned}[t]&\phantom{={}}0+10-1\\&=10-1\\&=9\end{aligned}\)\quad 收到10，再借出1给后一位。
\item 十分位传递退位：\(\begin{aligned}[t]&\phantom{={}}0+10-1\\&=10-1\\&=9\end{aligned}\)\quad 收到10，再借出1给后一位。
\item 百分位接收退位：\(\begin{aligned}[t]&\phantom{={}}0+10\\&=10\end{aligned}\)\quad 前一位的1变成本位的10。
\item 百分位：\(\begin{aligned}[t]&\phantom{={}}10-5\\&=5\end{aligned}\)\quad 本位写5。
\item 十分位：\(\begin{aligned}[t]&\phantom{={}}9-2\\&=7\end{aligned}\)\quad 本位写7。
\item 个位：\(\begin{aligned}[t]&\phantom{={}}9-3\\&=6\end{aligned}\)\quad 本位写6。
\item 十位：\(\begin{aligned}[t]&\phantom{={}}0-0\\&=0\end{aligned}\)\quad 本位写0。
\end{enumerate}
\par\smallskip
因此，\(\begin{aligned}[t]&\phantom{={}}10-3.25\\&=6.75\end{aligned}\)。
\par\medskip


\QHeading{计算过程：13. 解决问题}
\begin{enumerate}
\item 先由小明的身高求小红身高，再求小丽身高：
\Formula{\begin{aligned}[t]&\phantom{={}}1.45-0.08+0.12\\&=1.37+0.12\\&=1.49\text{（米）}\end{aligned}}
答：小丽身高1.49米。
\item 倒出的是一半油，桶重不变，所以前后总重之差是半桶油重。原来全部油重：
\Formula{\begin{aligned}[t]&\phantom{={}}(18.65-9.8)\times2\\&=8.85\times2\\&=17.7\text{（千克）}\end{aligned}}
桶重：\Formula{18.65-17.7=0.95\text{（千克）}.}答：桶重0.95千克。

答：原来桶里全部的油重17.7千克。
\end{enumerate}


}

\QHeading{第12题　用简便方法计算下面各题。}

\CalcBlock[方法一：简便计算]{1}{\begin{aligned}[t]&\phantom{={}}\color{black}2.35+1.48-0.35\\&\color{answerred}=2.35-0.35+1.48\\&\color{answerred}=(2.35-0.35)+1.48\\&\color{answerred}=2+1.48\\&\color{answerred}=3.48\end{aligned}}
\par
\MethodBlock{方法二：按运算顺序直接计算（对照）}{本法只作结果核对；本题要求简便计算，前面的简便解法仍保留。}{\begin{aligned}[t]&\phantom{={}}2.35+1.48-0.35\\
&=3.83-0.35\\
&=3.48\end{aligned}}


\CalcBlock[方法一：简便计算]{2}{\begin{aligned}[t]&\phantom{={}}\color{black}5.68+0.95+2.32\\&\color{answerred}=5.68+2.32+0.95\\&\color{answerred}=(5.68+2.32)+0.95\\&\color{answerred}=8+0.95\\&\color{answerred}=8.95\end{aligned}}
\MethodBlock{方法二：按运算顺序直接计算（对照）}{本法只作结果核对；本题要求简便计算，前面的简便解法仍保留。}{\begin{aligned}[t]&\phantom{={}}5.68+0.95+2.32\\
&=6.63+2.32\\
&=8.95\end{aligned}}

\CalcBlock[方法一：简便计算]{3}{\begin{aligned}[t]&\phantom{={}}\color{black}9.87-2.35-1.65\\&\color{answerred}=9.87-(2.35+1.65)\\&\color{answerred}=9.87-4\\&\color{answerred}=5.87\end{aligned}}
\MethodBlock{方法二：按运算顺序直接计算（对照）}{本法只作结果核对；本题要求简便计算，前面的简便解法仍保留。}{\begin{aligned}[t]&\phantom{={}}9.87-2.35-1.65\\
&=7.52-1.65\\
&=5.87\end{aligned}}

\CalcBlock[方法一：简便计算]{4}{\begin{aligned}[t]&\phantom{={}}\color{black}12.56-3.48-2.02\\&\color{answerred}=12.56-(3.48+2.02)\\&\color{answerred}=12.56-5.5\\&\color{answerred}=7.06\end{aligned}}
\MethodBlock{方法二：按运算顺序直接计算（对照）}{本法只作结果核对；本题要求简便计算，前面的简便解法仍保留。}{\begin{aligned}[t]&\phantom{={}}12.56-3.48-2.02\\
&=9.08-2.02\\
&=7.06\end{aligned}}

\CalcBlock[方法一：简便计算]{5}{\begin{aligned}[t]&\phantom{={}}\color{black}12.8+4.55+5.45\\&\color{answerred}=12.8+(4.55+5.45)\\&\color{answerred}=12.8+10\\&\color{answerred}=22.8\end{aligned}}
\MethodBlock{方法二：按运算顺序直接计算（对照）}{本法只作结果核对；本题要求简便计算，前面的简便解法仍保留。}{\begin{aligned}[t]&\phantom{={}}12.8+4.55+5.45\\
&=17.35+5.45\\
&=22.8\end{aligned}}

\CalcBlock[方法一：简便计算]{6}{\begin{aligned}[t]&\phantom{={}}\color{black}50.6-13.27-6.73\\&\color{answerred}=50.6-(13.27+6.73)\\&\color{answerred}=50.6-20\\&\color{answerred}=30.6\end{aligned}}
\MethodBlock{方法二：先处理括号或乘除，再多减再加回}{减去凑成的整数后，把多减的部分加回来。}{\begin{aligned}[t]&\phantom{={}}50.6-13.27-6.73\\
&=37.33-6.73\\
&=37.33-7+0.27\\
&=30.33+0.27\\
&=30.6\end{aligned}}

\par\noindent

\QHeading{第13题　解决问题。}

\Problem{① 小明身高1.45米，小红比小明矮0.08米，小丽比小红高0.12米。小丽身高多少米？}
\par


\MethodTitle{方法一：简便计算}
\Answer{{\Formula{\begin{aligned}[t]&\phantom{={}}1.45-0.08+0.12\\&=1.45+(0.12-0.08)\\&=1.45+0.04\\&=1.49\text{（米）}\end{aligned}}}\par 答：小丽身高1.49米。}

\MethodTitle{方法二：先求小丽比小明高多少}
\Answer{小丽相对小明先低0.08米，再高0.12米，合起来高0.04米。}
\Answer{\Formula{\begin{aligned}[t]&\phantom{={}}1.45+(0.12-0.08)\\
&=1.45+0.04\\
&=1.49\text{（米）}\end{aligned}}}
\Answer{答：小丽身高1.49米。}
\Problem{② 一桶油连桶重18.65千克，倒出一半油后，连桶重9.8千克。桶重多少千克？原来桶里全部的油重多少千克？}

\MethodTitle{方法一}


\MethodTitle{方法二：先由两次连桶重求桶重}
\Answer{把半桶油连桶重乘2，多出来的正好是一个桶，再求全部油重。}
\Answer{\textbf{桶重}\par \Formula{\begin{aligned}[t]&\phantom{={}}9.8\times 2-18.65\\
&=19.6-18.65\\
&=0.95\text{（千克）}\end{aligned}}}
\Answer{\textbf{全部油重}\par \Formula{\begin{aligned}[t]&\phantom{={}}18.65-0.95\\
&=17.7\text{（千克）}\end{aligned}}}
\Answer{答：桶重0.95千克。\par 答：原来全部的油重17.7千克。}

